%%=============================================================================
%% CAPITULO 4 -- APRESENTACAO, ANALISE E DISCUSSAO DOS DADOS
%%
%% Texto dummy (exercicio ia-6.1): trechos do pre-projeto do autor (versao 7).
%%=============================================================================

\section{Análise e discussão dos dados}

O protocolo TRTR será utilizado como referência: modelo treinado em dados reais
de treinamento e testado em dados reais de teste. O protocolo TSTR será o
experimento principal: modelo treinado em dados sintéticos e testado no mesmo
conjunto real de teste. Para classes desbalanceadas serão priorizados PR-AUC,
recall, F1 e, quando pertinente, intervalos de confiança; ROC-AUC será
complementar e não será utilizado isoladamente.

O artefato será considerado tecnicamente promissor se apresentar,
simultaneamente, evidências de fidelidade nas métricas selecionadas, utilidade
analítica adequada ao workload definido e ausência de sinais relevantes de
exposição segundo os critérios estabelecidos. Nenhuma dimensão será considerada
suficiente isoladamente. Um dataset que seja estatisticamente semelhante, mas não
preserve a tarefa analítica, não será classificado como útil; da mesma forma, um
dataset útil, mas com risco elevado de reprodução de registros, não será
considerado adequado para compartilhamento.
